% sample.tex — a showcase LaTeX document for testing gigvim's TeX support.
%
% This file exercises the pieces that texlab (LSP) and tex-fmt (formatter)
% care about: preamble/package loading, sectioning, math, floats (tables and
% figures), lists, cross-references, code listings and a bibliography.
%
% Build:
%   tectonic sample.tex            # single-shot, no bibtex dance (recommended)
%   latexmk -pdf sample.tex        # classic toolchain
%
% In gigvim: open this file, edit, and save — tex-fmt reformats on write
% (formatOnSave is enabled globally). :LspInfo should show texlab attached.

\documentclass[11pt,a4paper]{article}

\usepackage[utf8]{inputenc}
\usepackage[T1]{fontenc}
\usepackage{amsmath,amssymb}
\usepackage{graphicx}
\usepackage{booktabs}
\usepackage{listings}
\usepackage{xcolor}
\usepackage[colorlinks=true,linkcolor=blue,urlcolor=teal]{hyperref}

\lstset{
  basicstyle=\ttfamily\small,
  keywordstyle=\color{teal}\bfseries,
  commentstyle=\color{gray}\itshape,
  frame=single,
  breaklines=true,
  showstringspaces=false,
}

\title{A Sample \LaTeX{} Document}
\author{GigVim}
\date{\today}

\begin{document}

\maketitle

\begin{abstract}
  This document is a self-contained example used to verify that GigVim's
  TeX/LaTeX support works end to end: syntax highlighting via Treesitter,
  diagnostics and completion via \texttt{texlab}, and automatic formatting via
  \texttt{tex-fmt}. It touches most of the constructs you meet in a real
  paper---math, tables, figures, lists, code and references.
\end{abstract}

\tableofcontents

\section{Introduction}
\label{sec:intro}

Welcome to a minimal but representative \LaTeX{} document. We reference
Section~\ref{sec:math} for mathematics and Table~\ref{tab:results} for a small
results table. Inline math looks like $e^{i\pi} + 1 = 0$, and the editor should
highlight it distinctly from prose.

\section{Mathematics}
\label{sec:math}

Displayed equations are numbered and can be referenced. The Gaussian integral is
\begin{equation}
  \label{eq:gauss}
  \int_{-\infty}^{\infty} e^{-x^{2}} \, dx = \sqrt{\pi}.
\end{equation}

Aligned derivations keep their relation symbols lined up:
\begin{align}
  (a + b)^{2} &= a^{2} + 2ab + b^{2}, \\
  (a - b)^{2} &= a^{2} - 2ab + b^{2}.
\end{align}

Equation~\eqref{eq:gauss} is a staple of probability theory.

\section{Structure and Lists}
\label{sec:lists}

\subsection{Itemised}
\begin{itemize}
  \item Treesitter grammars: \texttt{latex} and \texttt{bibtex}.
  \item Language server: \texttt{texlab}.
  \item Formatter: \texttt{tex-fmt} (default) or \texttt{latexindent}.
\end{itemize}

\subsection{Enumerated}
\begin{enumerate}
  \item Open the file in GigVim.
  \item Edit some content.
  \item Save; watch \texttt{tex-fmt} tidy the source.
\end{enumerate}

\section{A Table}
\label{sec:table}

Table~\ref{tab:results} uses \texttt{booktabs} rules for a clean look.

\begin{table}[ht]
  \centering
  \begin{tabular}{lrr}
    \toprule
    Feature     & Enabled & Provider      \\
    \midrule
    Highlighting & yes     & treesitter    \\
    Diagnostics  & yes     & texlab        \\
    Formatting   & yes     & tex-fmt       \\
    \bottomrule
  \end{tabular}
  \caption{TeX capabilities wired up in GigVim.}
  \label{tab:results}
\end{table}

\section{A Figure}
\label{sec:figure}

Replace \texttt{example-image} with a real asset to render an actual graphic;
the \texttt{graphicx} package handles PNG, JPG and PDF inclusions.

\begin{figure}[ht]
  \centering
  % \includegraphics[width=0.6\linewidth]{example-image}
  \fbox{\parbox{0.6\linewidth}{\centering\vspace{2cm} figure placeholder \vspace{2cm}}}
  \caption{A placeholder standing in for an included image.}
  \label{fig:placeholder}
\end{figure}

\section{Code Listing}
\label{sec:code}

The Nix module that enables all of the above:

\begin{lstlisting}[language=Nix]
config.vim.languages.tex = {
  enable = true;
  format.enable = true;   # tex-fmt
  lsp.enable = true;      # texlab
  treesitter.enable = true;
};
\end{lstlisting}

\section{Conclusion}
\label{sec:conclusion}

If this compiles to a clean PDF and the editor shows highlighting, completion
and format-on-save, GigVim's LaTeX support is working. See
Section~\ref{sec:intro} to start over.

\end{document}
